\documentclass[conference]{IEEEtran}
\usepackage[utf8]{inputenc}
\usepackage[T1]{fontenc}
\usepackage{cite}
\usepackage{amsmath,amssymb}
\usepackage{graphicx}
\usepackage{booktabs}
\usepackage{url}

\title{Canary- and Rollback-Aware Composite Scoring for Predicting Operational Risk of Agentic NetOps Actions: An Artifact‑Backed Null and Reproducibility Report}

\author{
\IEEEauthorblockN{Autonomous AI Researcher}
\IEEEauthorblockA{CNSM 2026 Agentic AI Researcher Track}
}

\begin{document}

\maketitle

\begin{abstract}
We preregistered and executed a deterministic paired confirmatory study (CNSM2026-04-canary-rollback-metrics-v1) to test whether a linear composite (verifier\_coverage, canary\_sim\_pass, rollback\_success\_prob, LLM\_confidence, abstention\_utility) discriminates emulator‑measured outages better than any single submetric. Using a synthetic NetOps Task Bank and a deterministic validator, we ran 200 held‑out paired tasks (baseline vs guarded) and preserved exhaustive artifacts (raw responses, per‑episode scoring JSONs with validator\_trace, provider-call traces, execution manifest, and SHA‑256 hashes). Empirical result: all held‑out candidates validated successfully in both conditions (n11=200,n10=0,n01=0,n00=0). Paired‑bootstrap 95\% CI for the paired difference = [0.0,0.0] (bootstrap\_seed=1729,R=10000); McNemar two‑sided p = 1.0. The preregistered AUC‑ROC primary estimand is uncomputable on this held‑out set because labels are constant (no outage positives). Two preregistered artifacts required to recompute the composite‑vs‑submetric AUC pipeline (consolidated\_submetrics.csv and composite\_model\_weights.pkl) are absent from the preserved bundle. We (i) map principal claims to exact persisted artifact path(s) and SHA‑256 hash(es), (ii) present a deterministic reconstruction recipe limited to preserved artifacts, and (iii) provide an explicit artifact‑hygiene checklist (precise filenames and canonical fields) required to enable non‑degenerate discriminative evaluations in future runs.
\end{abstract}

\section{Introduction}
Agentic LLM-based NetOps promises to translate operator intent into executable configurations. Operational adoption requires demonstrable, reproducible evidence that the safety machinery around generation---deterministic validators, canary simulation, rollback trials, and abstention decision logic---actually predict runtime risk (outage probability) better than single-factor checks. Prior engineering and research emphasize generate--verify--repair and verifier orchestration; however, discriminative evaluations that connect per‑proposal submetrics to operational outages remain scarce. We preregistered CNSM2026-04-canary-rollback-metrics-v1 to test whether a linear composite score built from verifier\_coverage, canary\_sim\_pass, rollback\_success\_prob, LLM\_confidence, and abstention\_utility discriminates emulator outage (binary) better than any single submetric. The preregistration was machine‑frozen before execution (see Disclosure). Our contribution is an artifact‑grounded reproducibility report: a fully preregistered, deterministic run whose preserved artifacts substantiate a ceiling/null outcome on the held‑out test and reveal the exact persisted objects required (but missing) to complete the preregistered discriminative pipeline.

\section{Related Work}
We position this study with respect to verifier‑anchored GVR architectures and verifier orchestration. Representative verified prior works used for framing include: MeshAgent (constraint‑guided NetOps adaptation) --- DOI: 10.1145/3771567; GVR architectures for convergent correctness --- DOI: 10.2139/ssrn.6754899; sequential, cost‑aware verification (SEVER) --- DOI: 10.2139/ssrn.6810160; queueing/control for tool‑rich LLM serving --- DOI: 10.2139/ssrn.6438169; and a recent agentic NetOps survey used to motivate operational evaluation needs --- OpenAlex W7161354925. We relied on these verified sources for architectural framing only. The experimental evidence reported here is entirely artifact‑grounded in the archived run described below.

\section{Methodology}
Preregistration and frozen design. The preregistration JSON (SHA‑256 = 553509f5d61a7f0df674b676b1091ccbe734c2d142b52e65d81932fae8657cee) fixed: paired\_binary design; deterministic synthetic task bank split by content hash (train=300,val=100,test=200); planned\_episode\_count=400; shared\_initial\_generation (single shared candidate per pair); and maximum\_model\_calls=400. The composite feature vector and the analysis plan were prespecified: train L2 logistic regression (5‑fold CV) on train+val to obtain composite\_model\_weights.pkl; evaluate AUC‑ROC/AUPRC on the held‑out test; compare composite vs best single submetric using DeLong; use paired bootstrap (seed=1729,R=10000) to CI the paired difference in success rates. Execution. Adapter family: hosted\_netops\_gvr\_v1; declared model: gpt-5-mini (execution/model\_configuration.json SHA‑256 = 1acce92558edeb13cd97b75817329d332afe4a36d01d566f1a26c61b132923fa); model provider: openai\_responses; deterministic validator: deterministic\_netops\_validator\_v5. The run preserved raw responses (execution/responses/*.txt), per‑episode scoring JSONs (execution/scoring/*.json) containing validator\_trace.validation\_after, provider\_call traces (execution/provider\_calls/*.json), and a full execution manifest enumerating artifact SHA‑256s (execution/execution\_manifest.json). Scoring rule. Per the preregistration the validator returns per‑episode boolean valid plus parsed diagnostic fields (parsed\_command\_count, required\_command\_count, observed\_touched\_field\_count, normalized\_configuration) used to compute verifier\_coverage. Canary and rollback. The preregistration required per‑episode canary\_sim\_pass and rollback\_success\_prob to be persisted to reconstruct the composite; the run design included simulated canaries and scripted rollback trials, but their logs must be persisted under exact filenames listed below to enable recomputation.

\section{Results}
Execution integrity and retained artifacts. The execution manifest (execution/execution\_manifest.json as preserved) records planned\_episode\_count = 400, completed\_episode\_count = 400, failed\_episode\_count = 0, and model\_calls\_used = 200. Per‑episode validator outcomes are present as execution/scoring/<episode>.json files containing validator\_trace.validation\_after.normalized\_configuration and valid=true when the configuration satisfies the validator. Representative preserved per‑episode JSON: execution/scoring/task-000001-baseline.json (SHA‑256 = 90844c2baf63d1a1eb8ed5d6b26bfb38a7941ec4c4f1af9cb8930a33e9cc4478) --- see artifact\_examples in the evidence bundle. Paired confirmatory outcome (artifact‑grounded). The preserved paired contingency (analysis/paired\_contingency\_table.csv; SHA‑256 = 7bc1bd2a42b5ac3f7adb61db47cf8dbe0472f678032abcd1e7c44f92ea2de71c) is: n11 = 200, n10 = 0, n01 = 0, n00 = 0. analysis/condition\_summary.csv (SHA‑256 = 1cb072b4db7c8e29212ef28b97baccb66a507a7c8b8ceaad35b7e84f7976a951) records baseline\_success\_rate = 200/200 = 1.0 and guarded\_success\_rate = 200/200 = 1.0. Confirmatory statistics (artifact reproducible). From the preserved contingency we reproduce the paired difference estimate = 0.0; paired‑bootstrap percentile 95\% CI = [0.0,0.0] (bootstrap\_seed=1729,R=10000); McNemar two‑sided p = 1.0. These operations and their parameters are recorded in analysis/analysis\_log.jsonl (SHA‑256 = c56b737a55780f99fe36a9650ee69f18694a666348d4643c0dd8240e1757406d). Primary estimand status and missing persisted artifacts. The preregistered primary estimand (difference in AUC‑ROC: composite vs best single submetric) is undefined on the preserved held‑out test because labels are constant (no outage positives). Moreover, two preregistered objects required to deterministically reconstruct the composite‑vs‑submetric AUC pipeline were not persisted: consolidated\_submetrics.csv and composite\_model\_weights.pkl. Per‑submetric availability (artifact‑grounded). From preserved scoring and provider\_call JSONs we determine availability as follows: (1) verifier\_coverage --- PRESENT; extract from execution/scoring/<episode>.json \ensuremath{\rightarrow} validator\_trace.validation\_after.observed\_touched\_field\_count and required\_command\_count (e.g., execution/scoring/task-000001-baseline.json SHA‑256 above). (2) canary\_sim\_pass --- ABSENT: no per‑episode canary\_sim\_pass boolean persisted under the preregistered canary\_simulator\_log.json filename. (3) rollback\_success\_prob --- ABSENT: no rollback\_trial\_log.json with per‑episode rollback success probabilities is present. (4) LLM\_confidence --- ABSENT: execution/provider\_calls/*.json do not contain a numeric confidence field under preserved keys. (5) abstention\_utility --- ABSENT: no per‑episode abstention scoring artifact persisted. Deterministic reconstruction recipe (artifact‑limited). Using preserved execution/scoring/*.json and execution/responses/*.txt one can reproduce validator booleans and verifier\_coverage and thereby the paired contingency and McNemar/bootstrap results (already present). Reconstructing the composite AUC pipeline requires persisted consolidated\_submetrics.csv and composite\_model\_weights.pkl; absent these, the preregistered DeLong/AUC comparison cannot be computed from the bundle. Representative preserved artifacts (explicit paths + SHA‑256): execution/scoring/task-000001-baseline.json (90844c2b...), execution/responses/task-000001-shared-initial.txt (8f38c8be...), analysis/paired\_contingency\_table.csv (7bc1bd2a...), analysis/condition\_summary.csv (1cb072b4...), analysis/analysis\_log.jsonl (c56b737a...).

\section{Discussion}
Interpretation. The preserved ceiling/null outcome (all 200 held‑out pairs succeeded in both baseline and guarded conditions) limits inferential claims about discriminative predictive performance: ROC/AUPRC require positive labels and were therefore uncomputable on the archived test. This outcome highlights two distinct failure modes for confirmatory studies in this domain: (A) task-bank calibration producing insufficiently challenging instances (label imbalance leading to degenerate held‑out labels), and (B) omission of critical persisted submetric artifacts required to reconstruct composite discriminative analyses. Practical guidance and artifact‑hygiene checklist. To prevent repeat failure we supply an explicit, minimal set of exact filenames and canonical fields that preregistration must require and the execution pipeline must persist: 1) consolidated\_submetrics.csv --- columns: episode\_id,pair\_id,verifier\_coverage,canary\_sim\_pass,rollback\_success\_prob,LLM\_confidence,abstention\_utility. 2) composite\_model\_weights.pkl --- dict keys: feature\_names,scaling\_params\{mean,scale\},coefficients,intercept,training\_rng\_seed,cv\_spec. 3) canary\_simulator\_log.json --- per‑episode keys: episode\_id,canary\_seed,canary\_result\_boolean,canary\_trace\_path,canary\_trace\_sha256. 4) rollback\_trial\_log.json --- per‑episode keys: episode\_id,rollback\_seed,rollback\_success\_prob,rollback\_trace\_path,rollback\_trace\_sha256. 5) provider\_call confidence: ensure execution/provider\_calls/<id>.json includes a numeric confidence field named model\_confidence. 6) retain analysis/paired\_contingency\_table.csv and analysis/analysis\_log.jsonl (already present). Rationale: these exact filenames and canonical field names permit deterministic recomputation of composite weights, deterministic scoring of held‑out episodes, and ROC/AUPRC evaluation without ad hoc reconstruction. Threats to validity. Internal: shared\_initial candidate reuse reduces generation variance and can hide prompt sensitivity. Construct: synthetic emulator labels may not represent complex production outage semantics. External: single declared model and a single synthetic task bank limits generality. Conclusion. This artifact‑grounded null report is scientifically useful: it documents a fully preregistered deterministic run, makes every preserved artifact and its SHA‑256 explicit, reproduces the paired results exactly, identifies the precise persisted objects missing for the preregistered primary estimand, and supplies a compact checklist and deterministic recipe to enable non‑degenerate discriminative NetOps evaluations in future runs.

\section{Limitations}
\begin{itemize}
\item Synthetic task bank and deterministic emulator: tasks and validator may not capture real‑world nondeterminism, device heterogeneity, or latent failure modes.
\item Single declared model: results reflect gpt-5-mini under the declared adapter; generalization to other LLMs is untested.
\item Ceiling-effect null: zero observed outages in both conditions prevented discriminative evaluation of the preregistered AUC estimand.
\item Construct validity: emulator‑derived outage labels may not map one‑to‑one to production outage risk in real networks.
\item Artifact omission: consolidated\_submetrics.csv and composite\_model\_weights.pkl (required by preregistration) were not persisted, preventing recomputation of the preregistered composite discriminative analysis.
\end{itemize}

\section{Conclusion}
We preregistered and executed CNSM2026-04-canary-rollback-metrics-v1 and archived an exhaustive artifact set. The preserved evidence supports a concrete ceiling/null result on the held‑out test (n11=200,n10=0,n01=0,n00=0) with paired difference = 0.0, paired‑bootstrap 95\% CI = [0.0,0.0], and McNemar two‑sided p = 1.0. The preregistered AUC primary estimand could not be computed because (i) the held‑out labels are constant and (ii) two preregistered artifacts required to reconstruct the composite discriminative pipeline (consolidated\_submetrics.csv and composite\_model\_weights.pkl) were not persisted. We supply an exact artifact‑hygiene checklist and a deterministic reconstruction recipe (limited to preserved artifacts). Future confirmatory NetOps discriminative evaluations must (a) calibrate task banks to ensure label variation in held‑out splits, and (b) persist the exact per‑episode and training artifacts enumerated above. The scientific deliverable here is an artifact‑backed reproducibility report that maps claims to SHA‑256‑addressable evidence and documents precise operational lessons for repeatable agentic NetOps evaluation.

\section*{Disclosure Statement}
Human and autonomy boundary. A human Research Director selected the topic family (Generative AI / LLMs for NetOps) and approved the autonomy boundary prior to the autonomy lock. After the lock, all literature review, preregistration, experiment design freeze, code generation, experiment execution, artifact hashing, analysis, and manuscript drafting reported here were performed autonomously without manual human editing of technical content. Declared model and configuration. Declared model: gpt-5-mini (execution/model\_configuration.json SHA‑256 = 1acce92558edeb13cd97b75817329d332afe4a36d01d566f1a26c61b132923fa). Adapter family: hosted\_netops\_gvr\_v1. Deterministic validator: deterministic\_netops\_validator\_v5. Immutable initial master prompt. The immutable initial master prompt used by the autonomous researcher is archived at provenance/master\_prompt.txt (SHA‑256 = 1872df1e1805d2d96940456ca016bd665d1d5196add77f5acdf1582bb39b15ba) and was not modified after the autonomy lock. Preregistration. The preregistered study id is CNSM2026-04-canary-rollback-metrics-v1 and the preregistration JSON SHA‑256 is 553509f5d61a7f0df674b676b1091ccbe734c2d142b52e65d81932fae8657cee. Execution artifacts and preserved trace. The run preserved raw responses (execution/responses/*.txt), per‑episode scoring JSONs with validator\_trace (execution/scoring/*.json), provider\_call traces (execution/provider\_calls/*.json), and an execution manifest enumerating all artifacts and per‑artifact SHA‑256 (execution/execution\_manifest.json). Representative artifact paths and SHA‑256s used in the manuscript: execution/scoring/task-000001-baseline.json (SHA‑256 = 90844c2baf63d1a1eb8ed5d6b26bfb38a7941ec4c4f1af9cb8930a33e9cc4478), execution/responses/task-000001-shared-initial.txt (SHA‑256 = 8f38c8becd7e1e36d8902e9ba2f09218baad021d77365b20eb52914a4bdf4080), analysis/paired\_contingency\_table.csv (SHA‑256 = 7bc1bd2a42b5ac3f7adb61db47cf8dbe0472f678032abcd1e7c44f92ea2de71c), analysis/condition\_summary.csv (SHA‑256 = 1cb072b4db7c8e29212ef28b97baccb66a507a7c8b8ceaad35b7e84f7976a951), analysis/analysis\_log.jsonl (SHA‑256 = c56b737a55780f99fe36a9650ee69f18694a666348d4643c0dd8240e1757406d). Missing persisted artifacts and effect on primary estimand. The preregistration required consolidated\_submetrics.csv and composite\_model\_weights.pkl to be persisted to enable deterministic retraining and computation of the preregistered composite AUC comparison; neither file is present in the preserved manifest. Consequently, the preregistered primary estimand (difference in AUC‑ROC: composite vs best single submetric) could not be computed on the preserved test set. Safety and scope. All experiments were run on synthetic tasks and containerized emulators; no production devices were used. No technical restarts affected scientific outcomes (execution manifest: failed\_episode\_count = 0). The run preserved exhaustive logs and hash manifests to support independent artifact verification. Post‑lock human edits. No human manually edited technical manuscript content after the autonomy lock.

\begin{thebibliography}{99}
\bibitem{ref1} https://doi.org/10.1145/3771567
\bibitem{ref2} https://doi.org/10.2139/ssrn.6754899
\bibitem{ref3} https://doi.org/10.2139/ssrn.6810160
\bibitem{ref4} https://doi.org/10.2139/ssrn.6438169
\bibitem{ref5} Muhammad Bilal, Jon Crowcroft, Ruizhi Wang, Xiaolong Xu, Schahram Dustdar, "Large Language Models for Agentic NetOps and AIOps: Architectures, Evaluation, and Safety", 2026
\end{thebibliography}

\end{document}
